\section{Metodología de la segunda sesión}
El levantamiento se orientó por los principios de ingeniería de requerimientos y trazabilidad \cite{nuseibeh2000,wiegers2013}. Se inspeccionó estáticamente el código y esquema SQL depurados de PID-Original. Se diseñaron I-01 (entrevista semiestructurada de 16 preguntas) e I-02 (guía de observación de 14 ítems). I-02 se aplicó parcialmente a archivos remotos del repositorio; I-01 no se ha aplicado.

Se distinguieron evidencias de código (EV-COD), inspección (EV-OBS), aclaraciones históricas (CTX), propuestas (PROP) y simulaciones (SIM). Las afirmaciones negativas se limitan a los archivos revisados. Se diseñó un dataset completamente ficticio, SIM-001, para preparar escenarios de prueba, no para presentar resultados empíricos.

La calidad preliminar de los requerimientos se valoró mediante identificación, fuente, prioridad, aceptación verificable, consistencia, factibilidad y trazabilidad. La validación con actores académicos permanece pendiente.
